% ==============================================================================
% Chapter 9: MODULI PROBLEMS VIEWED OVER CYCLOTOMIC INTEGER RINGS
% ==============================================================================
\chapter{Moduli Problems Viewed over Cyclotomic Integer Rings}
\label{ch:09}

% Section 9.1
\section{Generalities}
\label{sec:9.1}

\kmitem\label{item:9.1.1}
Let $R$ be a ring, and let
\begin{equation}\label{eq:9.1.1.1}
\mathcal{F} : (\Sch/R) \mapsto (\Sets)
\end{equation}
be a contravariant functor which is representable by an affine $R$-scheme $\Spec(A)$. For example, $R = \Z$, and for any integer $N \ge 1$ the functor $\mu_N^\times$ ``primitive $N$'th roots of unity'' on $(\Sch/\Z)$:
\[
\begin{aligned}
\mu_N^\times(S) = {} & \text{the set of elements } \zeta \in \Gamma(S, \mathcal{O}_S) \\
& \text{such that } \Phi_N(\zeta) = 0
\end{aligned}
\]
where $\Phi_N(X) \in \Z[X]$ is the $N$'th cyclotomic polynomial. The corresponding $A$ is the ring
\[
\Z[\zeta_N] \overset{\mathrm{dfn}}{=} \Z[X]/(\Phi_N(X)) .
\]

\kmitem\label{item:9.1.2}
Now consider the moduli problem $[\mathcal{F}]$ on $(\mathrm{Ell}/R)$ defined by
\begin{equation}\label{eq:9.1.2.1}
[\mathcal{F}](E/S) \overset{\mathrm{dfn}}{=} \mathcal{F}(S) .
\end{equation}
It is relatively representable and affine over $(\mathrm{Ell}/R)$. Indeed for any $E/S/R$, we have
\begin{equation}\label{eq:9.1.2.2}
[\mathcal{F}]_{E/S} = S \tens{R} A .
\end{equation}
Obviously the moduli problem $[\mathcal{F}]$ is \emph{not} representable, simply because it is \emph{not} rigid: the group $\Aut(E/S)$ operates trivially on the set $[\mathcal{F}](E/S) = \mathcal{F}(S)$.

\begin{kprop}\label{prop:9.1.3}
There is a canonical morphism
\[
\Mcoarse([\mathcal{F}]) \to \Spec(A[j]) ,
\]
which is an isomorphism if $A$ is flat over $R$, or if $6 = 2 \times 3$ is invertible in $R$, or if $6$ is invertible in $A$.
\end{kprop}

\begin{proof}
The question is local on $R$, so we may assume some odd prime $\ell$ invertible in $R$. Then we have, with $G = \GL(2, \F_\ell)$,
\[
\begin{aligned}
\Mcoarse([\mathcal{F}]) & = \Mcoarse([\mathcal{F}], [\Gamma(\ell)])/G \\
& = (\Mcoarse([\Gamma(\ell)]) \tens{R} A)/G
\end{aligned}
\]
and the canonical map
\[
(\Mcoarse([\Gamma(\ell)]) \tens{R} A)/G \to (\Mcoarse([\Gamma(\ell)])/G) \tens{R} A = \Spec(R[j] \tens{R} A) .
\]
This last map is an isomorphism if $A$ is $R$-flat, or if the order of $G$ is invertible in $R$ or in $A$ (so if $A \ni 1/6$, take $\ell = 3$, giving $\#G = 48$).
\end{proof}

\kmitem\label{item:9.1.4}
Now consider a relatively representable and affine moduli problem $\mathcal{P}$ on $(\mathrm{Ell}/R)$, together with an $[\mathcal{F}]$-structure, i.e., a morphism of moduli problems on $(\mathrm{Ell}/R)$
\begin{equation}\label{eq:9.1.4.1}
\mathcal{P} \xrightarrow{\ \pi\ } [\mathcal{F}] ,
\end{equation}
i.e., a rule which to every level $\mathcal{P}$-structure on an $E/S/R$ attaches an $\mathcal{F}$-structure on $S$, compatibly with morphisms in $(\mathrm{Ell}/R)$. Because $\mathcal{P}$ is relatively representable, this means that for each $E/S/R$, we are given a morphism of $S$-schemes
\begin{equation}\label{eq:9.1.4.2}
\mathcal{P}_{E/S} \to S \tens{R} A ,
\end{equation}
i.e., for every representable problem $\mathcal{S}$ on $(\mathrm{Ell}/R)$, we have a morphism of $R$-schemes
\begin{equation}\label{eq:9.1.4.3}
\Mfrak(\mathcal{P}, \mathcal{S}) \to \Mfrak(\mathcal{S}) \tens{R} A .
\end{equation}
Passing to the map on coarse moduli schemes, we get a morphism of $R$-schemes
\begin{equation}\label{eq:9.1.4.4}
\Mcoarse(\mathcal{P}) \to \Mcoarse([\mathcal{F}]) ,
\end{equation}
which for $A/R$ flat or $1/6 \in R$ we may interpret as an $R$-morphism
\begin{equation}\label{eq:9.1.4.5}
\Mcoarse(\mathcal{P}) \to \Spec(A[j]) .
\end{equation}

\kmitem\label{item:9.1.5}
Now consider $(\mathrm{Ell}/A)$. An object of this category is an $E/S/A$. The given structure of $A$-scheme on $S$ provides (indeed ``is'') a canonical $\mathcal{F}$-structure on $S$ (e.g., if $A = \Z[X]/(\Phi_N(X)) = \Z[\zeta_N]$, then the inverse image of $\zeta_N \overset{\mathrm{dfn}}{=} X \bmod \Phi_N$ on any $\Z[\zeta_N]$-scheme $S$ is the canonical primitive $N$'th root of unity on $S$ we have in mind).

\kmitem\label{item:9.1.6}
Out of the given data
\[
\mathcal{P} \xrightarrow{\ \pi\ } [\mathcal{F}] ,
\]
we will now construct a new moduli problem, functorial in the original data, denoted
\begin{equation}\label{eq:9.1.6.1}
\mathcal{P}^{\mathrm{can}} \quad \text{on } (\mathrm{Ell}/A) ,
\end{equation}
as follows:
\begin{equation}\label{eq:9.1.6.2}
\begin{aligned}
\mathcal{P}^{\mathrm{can}}(E/S/A) = {} & \text{the set of } \mathcal{P}\text{-structures on } E/S \text{ which map} \\
& \text{by } \pi \text{ to the canonical } \mathcal{F}\text{-structure on } S .
\end{aligned}
\end{equation}
If we denote by $f_{\mathrm{taut}} \in \mathcal{F}(A)$ the tautological $\mathcal{F}$-structure on $A$ corresponding to the identity map in $\Hom_{R\text{-alg}}(A, A) = \mathcal{F}(A)$, then we can rewrite this as
\[
\mathcal{P}^{\mathrm{can}}(E/S/A) = \mathcal{P}(E/S) \underset{\mathcal{F}(S)}{\times} f_{\mathrm{taut}} .
\]
For example $[\mathcal{F}]^{\mathrm{can}}$ is the $[\Gamma(1)]$-moduli problem on $(\mathrm{Ell}/A)$.

\begin{kprop}\label{prop:9.1.7}
The moduli problem $\mathcal{P}^{\mathrm{can}}$ on $(\mathrm{Ell}/A)$ is relatively representable. It is representable if $\mathcal{P}$ is representable. For any representable moduli problem $\mathcal{S}$ on $(\mathrm{Ell}/R)$, we have a canonical isomorphism of schemes over $\Mfrak(\mathcal{S} \tens{R} A) = \Mfrak(\mathcal{S}) \tens{R} A$
\begin{equation*}
\begin{tikzcd}[row sep=2.2em, column sep=3.5em]
\Mfrak(\mathcal{P}^{\mathrm{can}}, \mathcal{S} \tens{R} A) \arrow[r, "\sim"] \arrow[d, "\mathrm{pr}_2"'] & \Mfrak(\mathcal{P}, \mathcal{S}) \arrow[d, "{(\mathrm{pr}_2, \pi)}"] \\
\Mfrak(\mathcal{S} \tens{R} A) \arrow[r, "\sim"] & \Mfrak(\mathcal{S}) \tens{R} A .
\end{tikzcd}
\end{equation*}
If $\mathcal{P}$ is affine over $(\mathrm{Ell}/R)$, then $\mathcal{P}^{\mathrm{can}}$ is affine over $(\mathrm{Ell}/A)$ and we have a canonical isomorphism of coarse moduli schemes as $A[j]$-schemes
\begin{equation*}
\begin{tikzcd}[row sep=2.0em, column sep=0.8em]
\Mcoarse(\mathcal{P}^{\mathrm{can}}) \arrow[rrrr, "\sim"] \arrow[rd] & & & & \Mcoarse(\mathcal{P}) \arrow[ld] \\
& \Mcoarse([\Gamma(1)] \tens{R} A) \arrow[rd, equal] & & \Mcoarse([\mathcal{F}]) \arrow[ld] & \\
& & \Spec(A[j]) . & &
\end{tikzcd}
\end{equation*}
\end{kprop}

\begin{proof}
As a moduli problem on $(\mathrm{Ell}/A)$, $\mathcal{P}^{\mathrm{can}}$ is a closed sub-problem of $\mathcal{P} \tens{R} A$. Over any $E/S/A$, $\mathcal{P}^{\mathrm{can}}$ is relatively represented by the fiber-product of the diagram
\begin{equation*}
\begin{tikzcd}[row sep=2.2em, column sep=3.5em]
\mathcal{P}_{E/S} \arrow[d] & \\
S \tens{R} A & S \arrow[l, "{(\mathrm{id}, \text{given})}"'] .
\end{tikzcd}
\end{equation*}
Now let $\mathcal{S}$ be a representable problem on $(\mathrm{Ell}/R)$, and $Z$ an $\Mfrak(\mathcal{S}) \tens{R} A$-scheme. Then $Z$ is provided with a structure of $A$-scheme, and $Z$ is provided with an elliptic curve $E/Z$ together with an $\mathcal{S}$-structure on $E/Z$. To map $Z$ as $\Mfrak(\mathcal{S}) \tens{R} A$-scheme to $\Mfrak(\mathcal{P}, \mathcal{S})$ is to give, independently, a level $\mathcal{P}$-structure on $E/Z$ which via $\pi$ lies over the tautological $\mathcal{F}$-structure on ``$Z$ as $A$-scheme.'' But this is exactly what it means to map $Z$ as $\Mfrak(\mathcal{S}) \tens{R} A = \Mfrak(\mathcal{S} \tens{R} A)$-scheme to $\Mfrak(\mathcal{P}^{\mathrm{can}}, \mathcal{S} \tens{R} A)$.

If $\mathcal{P}$ is representable, say by $E/T/R$, with universal element $\alpha \in \mathcal{P}(E/T)$, then $\pi(\alpha) \in \mathcal{F}(T) = \Hom_{R\text{-sch}}(T, \Spec(A))$ provides $T$ with the structure of $A$-scheme, say $\widetilde{T}/A$. Then on $E/\widetilde{T}/A$ the ``same'' element $\alpha$ lies in $\mathcal{P}^{\mathrm{can}}(E/\widetilde{T}/A)$, and the object $(E/\widetilde{T}/A, \alpha)$ visibly represents $\mathcal{P}^{\mathrm{can}}$.

If $\mathcal{P}$ is affine over $(\mathrm{Ell}/R)$, the fiber-product description of $\mathcal{P}^{\mathrm{can}}$ shows that $\mathcal{P}^{\mathrm{can}}$ is affine over $(\mathrm{Ell}/A)$. The isomorphism on coarse moduli schemes is constructed locally over $R$ by inverting an odd prime $\ell$, taking $\mathcal{S} = [\Gamma(\ell)]$, and then passing to $\GL(2, \F_\ell)$-invariants in the previous isomorphisms.
\end{proof}

\begin{kcor}\label{cor:9.1.8}
Let $Q$ be any property of schemes which is local for the etale topology. Then $\mathcal{P}$ has property $Q$ if and only if $\mathcal{P}^{\mathrm{can}}$ has property $Q$.
\end{kcor}

\begin{proof}
The problem is local on $R$, so we may assume an odd prime $\ell$ is invertible in $R$. Then $\mathcal{S} = [\Gamma(\ell)] \tens{\Z} R$ is etale and surjective over $(\mathrm{Ell}/A)$. Therefore $\mathcal{P}$ has property $Q$ if and only if the scheme $\Mfrak(\mathcal{P}, \mathcal{S})$ has property $Q$, while $\mathcal{P}^{\mathrm{can}}$ has property $Q$ if the scheme $\Mfrak(\mathcal{P}^{\mathrm{can}}, \mathcal{S} \tens{R} A)$ has it.
\end{proof}

\begin{kcor}\label{cor:9.1.9}
Let $P$ be any property of morphisms of schemes which is local for the etale topology on the base. Suppose that we have two relatively representable problems $\mathcal{P}, \mathcal{P}'$ on $(\mathrm{Ell}/R)$ both given with maps to $[\mathcal{F}]$, and a morphism of moduli problems on $(\mathrm{Ell}/R)$ respecting the $[\mathcal{F}]$-structures:
\begin{equation*}
\begin{tikzcd}[row sep=1.8em, column sep=1.5em]
\mathcal{P} \arrow[rr] \arrow[rd] & & \mathcal{P}' \arrow[ld] \\
& {[\mathcal{F}]} &
\end{tikzcd}
\end{equation*}
Then the induced map of relatively representable moduli problems on $(\mathrm{Ell}/A)$
\[
\mathcal{P}^{\mathrm{can}} \to (\mathcal{P}')^{\mathrm{can}}
\]
has property $P$ if and only if the original $\mathcal{P} \to \mathcal{P}'$ has property $P$. In particular, $\mathcal{P}^{\mathrm{can}}$ is of type $P$ over $(\mathrm{Ell}/A)$ if and only if the morphism $\mathcal{P} \to [\mathcal{F}]$ has property $P$.
\end{kcor}

\begin{proof}
Again the question is local on $R$, so we may assume an odd prime $\ell$ is invertible in $R$, and take $\mathcal{S} = [\Gamma(\ell)]$. Then $\mathcal{P} \to \mathcal{P}'$ has property $P$ if and only if the morphism of schemes
\[
\Mfrak(\mathcal{P}, \mathcal{S}) \to \Mfrak(\mathcal{P}', \mathcal{S})
\]
has property $P$, while $\mathcal{P}^{\mathrm{can}} \to (\mathcal{P}')^{\mathrm{can}}$ has property $P$ if and only if the same morphism of the same schemes
\[
\Mfrak(\mathcal{P}^{\mathrm{can}}, \mathcal{S} \tens{R} A) \to \Mfrak((\mathcal{P}')^{\mathrm{can}}, \mathcal{S} \tens{R} A)
\]
has property $P$.

In the special case $\mathcal{P}' = [\mathcal{F}]$, we have $[\mathcal{F}]^{\mathrm{can}} = {}$the moduli problem $[\Gamma(1)]$ on $(\mathrm{Ell}/A)$.
\end{proof}

\begin{kcor}\label{cor:9.1.10}
Suppose we are given $\mathcal{P} \xrightarrow{\ \pi\ } [\mathcal{F}]$ as above with $\mathcal{P}$ affine over $(\mathrm{Ell}/R)$, and a finite group $G$ acting trivially on $[\mathcal{F}]$ and acting on $\mathcal{P}$ by automorphisms which commute with $\pi$. Then $G$ acts on $\mathcal{P}^{\mathrm{can}}$ by functoriality, and we have a natural isomorphism of moduli problems on $(\mathrm{Ell}/A)$
\[
(\mathcal{P}/G)^{\mathrm{can}} \simeq \mathcal{P}^{\mathrm{can}}/G .
\]
\end{kcor}

\begin{proof}
The quotient $\mathcal{P}/G$ maps to $[\mathcal{F}]/G = [\mathcal{F}]$, so we have a $G$-equivariant commutative diagram
\begin{equation*}
\begin{tikzcd}[row sep=1.8em, column sep=1.5em]
\mathcal{P} \arrow[rr] \arrow[rd] & & \mathcal{P}/G \arrow[ld] \\
& {[\mathcal{F}]} &
\end{tikzcd}
\end{equation*}
Applying the $\mathcal{P} \mapsto \mathcal{P}^{\mathrm{can}}$ construction, we get by functoriality
\[
\mathcal{P}^{\mathrm{can}} \to (\mathcal{P}/G)^{\mathrm{can}} ,
\]
which is again $G$-equivariant, so defines
\[
\mathcal{P}^{\mathrm{can}}/G \to (\mathcal{P}/G)^{\mathrm{can}} .
\]
To show that this is an isomorphism, we may localize on $R$ and so assume some odd prime $\ell$ is invertible in $R$. Then $\mathcal{S} = [\Gamma(\ell)] \tens{\Z} R$ is etale and surjective over $(\mathrm{Ell}/R)$, and $\mathcal{S} \tens{R} A$ is etale and surjective over $(\mathrm{Ell}/A)$. It suffices to check that the map in question induces an isomorphism of schemes:
\[
\begin{tikzcd}[row sep=1.8em, column sep=2.5em, baseline=(current bounding box.south)]
\Mfrak(\mathcal{P}^{\mathrm{can}}/G, \mathcal{S} \tens{R} A) \arrow[r] \arrow[d, equal] & \Mfrak((\mathcal{P}/G)^{\mathrm{can}}, \mathcal{S} \tens{R} A) \arrow[d, equal] \\
\Mfrak(\mathcal{P}^{\mathrm{can}}, \mathcal{S} \tens{R} A)/G \arrow[d, equal] & \Mfrak(\mathcal{P}/G, \mathcal{S}) \arrow[d, equal] \\
\Mfrak(\mathcal{P}, \mathcal{S})/G \arrow[r, equal] & \Mfrak(\mathcal{P}, \mathcal{S})/G .
\end{tikzcd} \qedhere
\]
\end{proof}

\section{A descent situation}
\label{sec:9.2}

\kmitem\label{item:9.2.1}
Suppose that a finite group $G$ acts freely on $\mathcal{F}$, the representable functor on $(\Sch/R)$ represented by the $R$-algebra $A$, so that $A$ is an etale $G$-torsor over $A^G$. Then $G$ operates freely on $[\mathcal{F}]$, and the quotient $[\mathcal{F}]/G = [\mathcal{F}/G]$ where $\mathcal{F}/G$ is the representable functor on $(\Sch/R)$ represented by $A^G$.

Now suppose we are given a relatively representable and affine $\mathcal{P}$ on $\Ell/R$, together with an action of $G$ and a $G$-equivariant morphism of moduli problems on $\Ell/R$
\[
\mathcal{P} \xrightarrow{\ \pi\ } [\mathcal{F}] .
\]
Passing to the quotient by $G$, we get an induced map
\[
\mathcal{P}/G \to [\mathcal{F}/G] .
\]
Thus we can speak of $\mathcal{P}^{\mathrm{can}}$ on $\Ell/A$, and of $(\mathcal{P}/G)^{\mathrm{can}}$ on $\Ell/A^G$.

\begin{kprop}\label{prop:9.2.2}
In the above situation, we have a canonical isomorphism of moduli problems on $\Ell/A$
\[
(\mathcal{P}/G)^{\mathrm{can}} \tens{A^G} A \xrightarrow{\ \sim\ } \mathcal{P}^{\mathrm{can}} .
\]
\end{kprop}

\begin{proof}
Applying \eqref{prop:7.2.2}, we have a cartesian diagram of moduli problems on $\Ell/R$
\[
\begin{tikzcd}[row sep=1.8em, column sep=3em]
\mathcal{P} \arrow[r] \arrow[d] & \mathcal{P}/G \arrow[d] \\
{[\mathcal{F}]} \arrow[r] & {[\mathcal{F}/G]} ,
\end{tikzcd}
\]
which renders the assertion tautologous.
\end{proof}

\section{The situation near infinity}
\label{sec:9.3}

\begin{kprop}\label{prop:9.3.1}
Let $R$ be an excellent noetherian regular ring, $A$ a finite flat $R$-algebra which is itself regular, $\mathcal{F}$ the corresponding functor on $(\Sch/R)$ represented by $A$, and
\[
\mathcal{P} \xrightarrow{\ \pi\ } [\mathcal{F}]
\]
a relatively representable $\mathcal{P}$ on $\Ell/R$ together with an $[\mathcal{F}]$-structure. Suppose that $\mathcal{P}$ is finite over $\Ell/R$, and normal near infinity. Then $\mathcal{P}^{\mathrm{can}}$ is finite over $\Ell/R$, and it is normal near infinity. The isomorphism of coarse moduli schemes
\[
\begin{tikzcd}[row sep=1.8em, column sep=2.5em]
\Mcoarse(\mathcal{P}^{\mathrm{can}}) \arrow[rr, "\sim"] \arrow[rd] & & \Mcoarse(\mathcal{P}) \arrow[dd] \\
& \Spec(A[j]) \arrow[rd] & \\
& & \Spec(R[j])
\end{tikzcd}
\]
extends to an isomorphism of compactified coarse moduli schemes
\[
\begin{tikzcd}[row sep=1.8em, column sep=2.5em]
\overline{\Mcoarse}(\mathcal{P}^{\mathrm{can}}) \arrow[rr, "\sim"] \arrow[rd] & & \overline{\Mcoarse}(\mathcal{P}) \arrow[ld] \arrow[dd] \\
& \PP^1_A \arrow[rd] & \\
& & \PP^1_R ,
\end{tikzcd}
\]
and induces isomorphisms of $A[[q]]$-schemes
\[
\begin{tikzcd}[row sep=1.8em, column sep=2.5em]
\widehat{\mathrm{Cusps}}(\mathcal{P}^{\mathrm{can}}) \arrow[rr, "\sim"] \arrow[rd] & & \widehat{\mathrm{Cusps}}(\mathcal{P}) \arrow[ld] \arrow[dd] \\
& \Spec(A[[q]]) \arrow[rd] & \\
& & \Spec(R[[q]]) ,
\end{tikzcd}
\]
isomorphisms of $A((q))$-schemes
\[
\begin{tikzcd}[row sep=1.8em, column sep=2.5em]
\widehat{\mathrm{Cusps}}(\mathcal{P}^{\mathrm{can}})_{A((q))} \arrow[rr, "\sim"] \arrow[rd] & & \widehat{\mathrm{Cusps}}(\mathcal{P})_{R((q))} \arrow[ld] \arrow[dd] \\
& \Spec(A((q))) \arrow[rd] & \\
& & \Spec(R((q))) ,
\end{tikzcd}
\]
and isomorphisms of $A$-schemes
\[
\begin{tikzcd}[row sep=1.8em, column sep=2.5em]
\mathrm{Cusps}(\mathcal{P}^{\mathrm{can}}) \arrow[rr, "\sim"] \arrow[rd] & & \mathrm{Cusps}(\mathcal{P}) \arrow[ld] \arrow[dd] \\
& \Spec(A) \arrow[rd] & \\
& & \Spec(R) .
\end{tikzcd}
\]
\end{kprop}

\begin{proof}
Because $A$ is finite over $R$ and normal, $\PP^1_A$ is normal, and finite over $\PP^1_R$. Therefore it is the same to normalize a neighborhood of infinity $\overline{U} \subset \PP^1_R$ in the finite $A[j]$-scheme $\Mcoarse(\mathcal{P})$, or to normalize $\overline{U} \tens{R} A \subset \PP^1_A$ in the same scheme. This gives the required isomorphism
\[
\begin{tikzcd}[row sep=1.8em, column sep=2.5em]
\overline{\Mcoarse}(\mathcal{P}^{\mathrm{can}}) \arrow[rr, "\sim"] \arrow[rd] & & \overline{\Mcoarse}(\mathcal{P}) \arrow[ld] \arrow[dd] \\
& \PP^1_A \arrow[rd] & \\
& & \PP^1_R .
\end{tikzcd}
\]
The scheme of cusps is the reduced subscheme ``$1/j = 0$'', which gives
\[
\begin{tikzcd}[row sep=1.8em, column sep=2.5em]
\mathrm{Cusps}(\mathcal{P}^{\mathrm{can}}) \arrow[rr, "\sim"] \arrow[rd] & & \mathrm{Cusps}(\mathcal{P}) \arrow[ld] \arrow[dd] \\
& \Spec(A) \arrow[rd] & \\
& & \Spec(R) .
\end{tikzcd}
\]
Completing along this subscheme gives
\[
\begin{tikzcd}[row sep=1.8em, column sep=2.5em]
\widehat{\mathrm{Cusps}}(\mathcal{P}^{\mathrm{can}}) \arrow[rr, "\sim"] \arrow[rd] & & \widehat{\mathrm{Cusps}}(\mathcal{P}) \arrow[ld] \arrow[dd] \\
& \Spec(A[[q]]) \arrow[rd] & \\
& & \Spec(R[[q]]) ,
\end{tikzcd}
\]
and inverting $q$ then gives
\[
\begin{tikzcd}[row sep=1.8em, column sep=2.5em]
\widehat{\mathrm{Cusps}}(\mathcal{P}^{\mathrm{can}})_{A((q))} \arrow[rr, "\sim"] \arrow[rd] & & \widehat{\mathrm{Cusps}}(\mathcal{P})_{R((q))} \arrow[ld] \arrow[dd] \\
& \Spec(A((q))) \arrow[rd] & \\
& & \Spec(R((q))) .
\end{tikzcd} \qedhere
\]
\end{proof}

\begin{kprop}\label{prop:9.3.2}
Let $R$ be an arbitrary ring, $A$ a finite free $R$-algebra, $\mathcal{F}$ the functor on $(\Sch/A)$ represented by $A$, and
\[
\mathcal{P} \xrightarrow{\ \pi\ } [\mathcal{F}]
\]
an $[\mathcal{F}]$-structure on a relatively representable moduli problem $\mathcal{P}$ on $\Ell/R$. Then we have an isomorphism of $A((q))$-schemes
\[
\mathcal{P}^{\mathrm{can}}_{\Tate(q)/A((q))} \xrightarrow{\ \sim\ } \mathcal{P}_{\Tate(q)/R((q))} .
\]
\end{kprop}

\begin{proof}
This is just \eqref{prop:9.1.7} applied to $\mathcal{S} = {}$the problem represented by $\Tate(q)/R((q))$. Because $A$ is a free $R$-module of finite rank, we have
\[
R((q)) \tens{R} A \xrightarrow{\ \sim\ } A((q)) ,
\]
so that $\mathcal{S} \tens{R} A$ is indeed the problem represented by $\Tate(q)/A((q))$.
\end{proof}


\section{Applications to the basic moduli problems}
\label{sec:9.4}

\kmitem\label{item:9.4.1}
Let $N \ge 1$ be an integer. Consider the finite free $\Z$-algebra
\[
A = \Z[\zeta_N] \xlongequal{\mathrm{dfn}} \Z[X]/(\Phi_N(X)) ,
\]
which represents the functor $\mu_N^\times$ of ``primitive $N$'th roots of unity'' on $(\Sch)$. The $e_N$-pairing defines a morphism of moduli problems on $\Ell$
\[
[\Gamma(N)] \to \mu_N^\times
\]
\[
E/S, (P,Q) \mapsto e_N(P,Q) \in \mu_N^\times(S) .
\]
The group $\GL(2, \Z/N\Z)$ acts on the right on $[\Gamma(N)]$; $g = \begin{pmatrix} a & b \\ c & d \end{pmatrix}$ acts on the Drinfeld $N$-basis $(P,Q)$ by
\[
(P,Q) \mapsto (P,Q)\begin{pmatrix} a & b \\ c & d \end{pmatrix} = (P_g, Q_g) .
\]
In terms of the corresponding homomorphism
\[
\phi \colon (\Z/N\Z)^2 \to E[N](S)
\]
\[
\phi \begin{pmatrix} x \\ y \end{pmatrix} = xP + yQ = (P,Q)\begin{pmatrix} x \\ y \end{pmatrix} ,
\]
the action is
\[
\phi \mapsto \phi \circ g .
\]
Because $e_N$ is alternating, we have
\[
e_N(P_g, Q_g) = e_N(aP + cQ, bP + dQ) = (e_N(P,Q))^{\det(g)} .
\]

\kmitem\label{item:9.4.2}
Let $G \subset \GL(2,\Z/N\Z)$ be a subgroup, and $\det(G) \subset (\Z/N\Z)^\times$ the image of $G$ under the determinant. Let $(\Z[\zeta_N])^{\det(G)}$ denote the subring of invariants of $\det(G)$ acting on $\Z[\zeta_N]$. The ring inclusion
\[
(\Z[\zeta_N])^{\det(G)} \hookrightarrow \Z[\zeta_N] ,
\]
defines a morphism of $\Z$-schemes
\[
\begin{tikzcd}[row sep=1.8em, column sep=2.5em]
\mu_N^\times = \Spec(\Z[\zeta_N]) \arrow[r] & \mu_N^\times/\det(G) \arrow[d, equal, "\mathrm{dfn}"] \\
& \Spec(\Z[\zeta_N])/\det G \arrow[d, equal] \\
& \Spec((\Z[\zeta_N])^{\det(G)})
\end{tikzcd}
\]
which by composition with
\[
[\Gamma(N)] \to [\mu_N^\times]
\]
defines
\[
[\Gamma(N)] \to [\mu_N^\times/\det(G)] .
\]
This map is respected by the action of $G$, so we obtain
\[
[\Gamma(N)]/G \to [\mu_N^\times/\det(G)] .
\]
Applying the $\mathcal{P} \mapsto \mathcal{P}^{\mathrm{can}}$ construction to this situation, we obtain a relatively representable moduli problem
\[
([\Gamma(N)]/G)^{\mathrm{can}} \quad\text{on } \Ell/(\Z[\zeta_N])^{\det(G)} .
\]

\kmitem\label{item:9.4.3}
In all those cases where we know a direct modular description of a quotient $[\Gamma(N)]/G$ of this sort, we can give a modular description of the morphism
\[
[\Gamma(N)]/G \to [\mu_N^\times/\det(G)] ,
\]
and consequently we can give a directly modular description of the moduli problem $([\Gamma(N)]/G)^{\mathrm{can}}$ on $\Ell/(\Z[\zeta_N])^{\det(G)}$. Here is a short list of such descriptions:

\par\addvspace{6pt plus 2pt minus 1pt}\noindent\phantomsection\label{item:9.4.3.1}(9.4.3.1)\enspace $N \ge 1$ arbitrary, $G$ trivial. $[\Gamma(N)]^{\mathrm{can}}$ is the moduli problem on $\Ell/\Z[\zeta_N]$ defined by
\[
\begin{aligned}
E/S/\Z[\zeta_N] \mapsto {} & \text{Drinfeld $N$-bases } (P,Q) \text{ on} \\
& E/S \text{ with } e_N(P,Q) = \zeta_N .
\end{aligned}
\]

\par\addvspace{6pt plus 2pt minus 1pt}\noindent\phantomsection\label{item:9.4.3.2}(9.4.3.2)\enspace $N \ge 1$ arbitrary, $G = \begin{pmatrix} 1 & * \\ 0 & 1 \end{pmatrix}$ the upper unipotent subgroup. The quotient $[\Gamma(N)]/G$ is $[\mathrm{bal.}\ \Gamma_1(N)]$, and $[\mathrm{bal.}\ \Gamma_1(N)]^{\mathrm{can}}$ is the moduli problem on $\Ell/\Z[\zeta_N]$ defined by
\[
\begin{aligned}
E/S/\Z[\zeta_N] \mapsto {} & \text{dual cyclic $N$-isogenies over $S$} \\
& P;\ E \kmdual{\pi}{\pi^t} E';\ P' \text{ together with} \\
& \text{generators $P$ of $\Ker\pi$, $P'$ of} \\
& \Ker\pi^t, \text{ whose canonical pairing} \\
& \text{is } \langle P, P' \rangle = \zeta_N .
\end{aligned}
\]

\par\addvspace{6pt plus 2pt minus 1pt}\noindent\phantomsection\label{item:9.4.3.3}(9.4.3.3)\enspace $N \ge 1$ arbitrary, $G = \begin{pmatrix} 1 & * \\ 0 & * \end{pmatrix}$. The quotient $[\Gamma(N)]/G$ is $[\Gamma_1(N)]$, but $\det(G) = (\Z/N\Z)^\times$, so $(\Z[\zeta_N])^{\det G} = \Z$, and $[\Gamma_1(N)]^{\mathrm{can}}$ is just $[\Gamma_1(N)]$, on $\Ell/\Z$.

\par\addvspace{6pt plus 2pt minus 1pt}\noindent\phantomsection\label{item:9.4.3.4}(9.4.3.4)\enspace $N \ge 1$ arbitrary, $G = \begin{pmatrix} * & * \\ 0 & * \end{pmatrix}$. The quotient is $[\Gamma_0(N)]$, but again $\det(G) = (\Z/N\Z)^\times$, so $[\Gamma_0(N)]^{\mathrm{can}}$ is just $[\Gamma_0(N)]$, on $\Ell/\Z$.

\par\addvspace{6pt plus 2pt minus 1pt}\noindent\phantomsection\label{item:9.4.3.5}(9.4.3.5)\enspace $N \ge 1$ arbitrary, $d$ a divisor of $N$, $G$ the subgroup of all matrices congruent $\bmod\ d$ to the identity. Then $\det(G)$ is the subgroup $(1+(d)) \cap (\Z/N\Z)^\times$, the quotient $[\Gamma(N)]/G$ is $[\Gamma(d)]$, and $[\Gamma(d)]^{\mathrm{can}}$ is the moduli problem on $\Ell/\Z[\zeta_d]$ described in (\ref{item:9.4.3.1}) above, with $N$ replaced by $d$.

\par\addvspace{6pt plus 2pt minus 1pt}\noindent\phantomsection\label{item:9.4.3.6}(9.4.3.6)\enspace $N = p^n$, $a,b$ two integers in $[0,n]$, $G_{a,b}$ the subgroup of all invertible matrices of the form
\[
\begin{pmatrix} 1 + (p^a) & * \\ 0 & 1 + (p^b) \end{pmatrix} .
\]
Then $\det(G)$ is the subgroup $1 + (p^{\min(a,b)})$ of $(\Z/p^n\Z)^\times$, and $[\Gamma(p^n)]/G_{a,b}$ is $[\Gamma_0(p^n); a,b]$. The problem $[\Gamma_0(p^n); a,b]^{\mathrm{can}}$ on $\Ell/\Z[\zeta_{p^{\min(a,b)}}]$ is
\[
\begin{aligned}
E/S/\Z[\zeta_{p^{\min(a,b)}}] \mapsto {} & \text{dual cyclic $p^n$ isogenies together with} \\
& \text{generators of the last $a$-step (resp.\ $b$-step)} \\
& \text{parts of their standard factorization}
\end{aligned}
\]
\[
E = E_0 \kmdual{\pi_{0,1}}{\lambda_{1,0}} \cdots \kmdual{\pi_{n-1,n}}{\lambda_{n,n-1}} E_n
\]
\[
\begin{aligned}
P_{n-a} & : \text{generator of } \Ker(\pi_{n-a,n}) \\
Q_b & : \text{generator of } \Ker(\lambda_{b,0}) ,
\end{aligned}
\]
such that
\begin{align*}
\text{if } a \le b, \quad \zeta_{p^a} & = \langle P_{n-a}, p^{a-b}\pi_{b,n} Q_b \rangle_{\pi_{n-a,n}} \\
\text{if } b \le a, \quad \zeta_{p^b} & = \langle p^{a-b}\lambda_{n-a,0} P_{n-a}, Q_b \rangle_{\pi_{0,b}} .
\end{align*}

\par\addvspace{6pt plus 2pt minus 1pt}\noindent\phantomsection\label{item:9.4.3.7}(9.4.3.7)\enspace $N_1, N_2$ two relatively prime integers, $G_1 \subset \GL(2,\Z/N_1\Z)$ and $G_2 \subset \GL(2,\Z/N_2\Z)$ subgroups such that the quotients $\mathcal{P}_1 = [\Gamma(N_1)]/G_1$ and $\mathcal{P}_2 = [\Gamma(N_2)]/G_2$ are both flat over $\Ell$. Then the simultaneous problem $\mathcal{P} = (\mathcal{P}_1, \mathcal{P}_2)$ is the quotient of $[\Gamma(N_1 N_2)]$ by the product subgroup $G_1 \times G_2$ of $\GL(2,\Z/N_1 N_2\Z)$, the ring $A = (\Z[\zeta_N])^{\det G}$ is the tensor product over $\Z$ of the rings $A_i = (\Z[\zeta_{N_i}])^{\det(G_i)}$ for $i = 1,2$, and $\mathcal{P}^{\mathrm{can}}$ on $\Ell/A$ is the simultaneous problem
\[
(\mathcal{P}_1^{\mathrm{can}} \tens{A_1} A,\ \mathcal{P}_2^{\mathrm{can}} \tens{A_2} A) .
\]
